\documentclass[a4paper,12pt]{ctexart}
\usepackage{../teaching-style}

\begin{document}
\pagestyle{empty}

\maintitle{计算基本功：四则运算}
\begin{center}
  \large 适用于4\~{}5年级 \quad 纯练习册
\end{center}
\vspace{0.5cm}

\chaptertitle{第一单元\quad 整数加减简算}

\sectiontitle{方法回顾}
\knowbox{凑整法：把相加能凑成整十、整百、整千的数先加。\\
基准数法：几个接近的数相加，取基准数再调整。\\
加法交换律：$a+b=b+a$ \qquad 加法结合律：$(a+b)+c=a+(b+c)$}

\bigskip

\begin{exercise}
\item 用凑整法计算：
  \begin{subexercise}
    \item $29+35+11$
    \item $48+57+52$
    \item $136+89+64$
    \item $275+347+25$
    \item $1234+567+8766$
    \item $888+999+112$
  \end{subexercise}

\item 用凑整法计算：
  \begin{subexercise}
    \item $199+299+99$
    \item $502+398+605$
    \item $1999+299+39$
    \item $9998+998+98$
    \item $789+997+88$
    \item $3456+1999+544$
  \end{subexercise}

\item 用基准数法计算：
  \begin{subexercise}
    \item $52+49+51+50+48$
    \item $97+102+98+101+99$
    \item $203+198+205+197+200$
    \item $302+299+305+298+301$
  \end{subexercise}

\item 计算：
  \begin{subexercise}
    \item $735-199$
    \item $864-398$
    \item $1200-507$
    \item $2000-345-655$
    \item $1563-299-401$
    \item $2458-367-633$
  \end{subexercise}

\item 计算：
  \begin{subexercise}
    \item $826-79-21$
    \item $574-136-64$
    \item $932-287-113$
    \item $745-98-202$
    \item $1685-347-253$
    \item $4321-567-433$
  \end{subexercise}

\item 计算：
  \begin{subexercise}
    \item $123+456-23$
    \item $789+234-89$
    \item $567+345-167$
    \item $901+234-201$
    \item $678+432-278$
    \item $345+567-145$
  \end{subexercise}

\item 计算：
  \begin{subexercise}
    \item $173-56+27$
    \item $284-75+16$
    \item $365-89-65$
    \item $452-67-52$
    \item $567-88-67+88$
    \item $789-123+211-77$
  \end{subexercise}
\end{exercise}

\newpage

\chaptertitle{第二单元\quad 整数乘除简算}

\sectiontitle{方法回顾}
\knowbox{乘法交换律：$a\times b=b\times a$ \qquad 乘法结合律：$(a\times b)\times c=a\times(b\times c)$\\
乘法分配律：$a\times(b+c)=a\times b+a\times c$ \qquad $a\times(b-c)=a\times b-a\times c$\\
商不变性质：被除数和除数同时乘或除以同一个非零数，商不变。}

\bigskip

\begin{exercise}
\item 用乘法结合律简算：
  \begin{subexercise}
    \item $25\times17\times4$
    \item $125\times23\times8$
    \item $5\times37\times20$
    \item $25\times13\times4$
    \item $125\times7\times8$
    \item $4\times9\times25$
  \end{subexercise}

\item 用乘法分配律简算：
  \begin{subexercise}
    \item $36\times99$
    \item $48\times101$
    \item $25\times199$
    \item $125\times801$
    \item $54\times102$
    \item $37\times98$
  \end{subexercise}

\item 用乘法分配律简算：
  \begin{subexercise}
    \item $25\times(40+4)$
    \item $125\times(80+8)$
    \item $45\times(100-2)$
    \item $64\times(50-1)$
    \item $36\times34+36\times66$
    \item $47\times58+47\times42$
  \end{subexercise}

\item 简算：
  \begin{subexercise}
    \item $56\times99+56$
    \item $72\times101-72$
    \item $38\times99+38$
    \item $95\times101-95$
    \item $88\times125$
    \item $32\times25$
  \end{subexercise}

\item 用商不变性质简算：
  \begin{subexercise}
    \item $400\div25$
    \item $800\div125$
    \item $3000\div125$
    \item $6000\div25$
    \item $4200\div(42\times25)$
    \item $3600\div(36\times125)$
  \end{subexercise}

\item 计算：
  \begin{subexercise}
    \item $1250\div25\div5$
    \item $1800\div(18\times25)$
    \item $360\div45$
    \item $480\div32$
    \item $540\div36$
    \item $720\div48$
  \end{subexercise}

\item 计算：
  \begin{subexercise}
    \item $(125+50)\times8$
    \item $25\times(40-8)$
    \item $125\times32\times25$
    \item $99\times99+99$
    \item $101\times87-87$
    \item $98\times34+2\times34$
  \end{subexercise}
\end{exercise}

\newpage

\chaptertitle{第三单元\quad 小数四则混合}

\sectiontitle{方法回顾}
\knowbox{小数加减法：小数点对齐，按位加减。\\
小数乘法：先按整数乘，再点小数点（因数小数位数之和）。\\
小数除法：除数变整数，被除数同步扩大。}

\bigskip

\begin{exercise}
\item 计算：
  \begin{subexercise}
    \item $3.6+4.8$
    \item $7.2-3.9$
    \item $12.5+8.75$
    \item $23.6-7.84$
    \item $5.67+3.45$
    \item $10.2-6.78$
  \end{subexercise}

\item 计算：
  \begin{subexercise}
    \item $3.5\times4$
    \item $2.5\times1.2$
    \item $0.25\times4.8$
    \item $1.25\times0.8$
    \item $3.6\times0.5$
    \item $0.75\times0.4$
  \end{subexercise}

\item 计算：
  \begin{subexercise}
    \item $7.2\div0.8$
    \item $4.5\div0.15$
    \item $12.6\div0.3$
    \item $3.24\div1.2$
    \item $0.81\div0.27$
    \item $5.6\div0.35$
  \end{subexercise}

\item 简算：
  \begin{subexercise}
    \item $2.5\times3.7\times4$
    \item $1.25\times0.7\times8$
    \item $0.25\times4.4$
    \item $0.125\times8.8$
    \item $3.6\times99+3.6$
    \item $4.8\times101-4.8$
  \end{subexercise}

\item 简算：
  \begin{subexercise}
    \item $5.6\times3.2+5.6\times6.8$
    \item $7.2\times4.5+2.8\times4.5$
    \item $12.5\times3.2\times0.25$
    \item $2.5\times(40+0.4)$
    \item $9.9\times5.6$
    \item $10.1\times3.4$
  \end{subexercise}

\item 计算：
  \begin{subexercise}
    \item $3.6+4.8-2.5$
    \item $12.5-3.6+7.5$
    \item $7.2\times0.5\div0.3$
    \item $4.5\div0.9\times1.2$
    \item $(5.6+4.4)\times0.25$
    \item $(8.8-3.8)\div0.5$
  \end{subexercise}

\item 计算：
  \begin{subexercise}
    \item $0.25\times(4+0.8)$
    \item $1.25\times(0.8-0.4)$
    \item $4.8\div0.6+3.2$
    \item $7.5-2.5\times0.4$
    \item $12.6\div(0.2\times0.7)$
    \item $3.6\times(0.5+0.25)$
  \end{subexercise}
\end{exercise}

\newpage

\chaptertitle{第四单元\quad 分数加减}

\sectiontitle{方法回顾}
\knowbox{同分母分数相加减：分母不变，分子相加减。\\
异分母分数相加减：先通分，再按同分母法则计算。\\
带分数加减：整数部分和分数部分分别加减。\\
加法运算律对分数同样适用。}

\bigskip

\begin{exercise}
\item 计算：
  \begin{subexercise}
    \item $\dfrac{2}{5}+\dfrac{1}{5}$
    \item $\dfrac{7}{9}-\dfrac{2}{9}$
    \item $\dfrac{3}{8}+\dfrac{5}{8}$
    \item $\dfrac{11}{12}-\dfrac{5}{12}$
    \item $\dfrac{5}{7}+\dfrac{4}{7}$
    \item $\dfrac{13}{15}-\dfrac{8}{15}$
  \end{subexercise}

\item 计算：
  \begin{subexercise}
    \item $\dfrac{1}{2}+\dfrac{1}{3}$
    \item $\dfrac{2}{3}-\dfrac{1}{4}$
    \item $\dfrac{3}{4}+\dfrac{1}{6}$
    \item $\dfrac{5}{6}-\dfrac{3}{8}$
    \item $\dfrac{2}{5}+\dfrac{3}{10}$
    \item $\dfrac{7}{12}-\dfrac{1}{3}$
  \end{subexercise}

\item 计算：
  \begin{subexercise}
    \item $\dfrac{5}{12}+\dfrac{7}{18}$
    \item $\dfrac{9}{20}-\dfrac{3}{15}$
    \item $\dfrac{11}{24}+\dfrac{5}{36}$
    \item $\dfrac{13}{30}-\dfrac{7}{45}$
    \item $\dfrac{8}{21}+\dfrac{5}{14}$
    \item $\dfrac{17}{28}-\dfrac{9}{35}$
  \end{subexercise}

\item 计算：
  \begin{subexercise}
    \item $2\dfrac{1}{3}+1\dfrac{1}{4}$
    \item $3\dfrac{2}{5}-1\dfrac{1}{3}$
    \item $5\dfrac{3}{8}+2\dfrac{1}{6}$
    \item $4\dfrac{5}{12}-2\dfrac{1}{4}$
    \item $6\dfrac{2}{7}+3\dfrac{1}{2}$
    \item $8\dfrac{3}{10}-5\dfrac{2}{15}$
  \end{subexercise}

\item 简算：
  \begin{subexercise}
    \item $\dfrac{3}{7}+\dfrac{2}{5}+\dfrac{4}{7}$
    \item $\dfrac{5}{9}+\dfrac{3}{8}+\dfrac{4}{9}+\dfrac{5}{8}$
    \item $\dfrac{7}{12}+\dfrac{5}{11}+\dfrac{5}{12}$
    \item $\dfrac{9}{13}+\dfrac{7}{10}-\dfrac{9}{13}$
  \end{subexercise}

\item 计算：
  \begin{subexercise}
    \item $\dfrac{1}{2}+\dfrac{1}{4}+\dfrac{1}{8}$
    \item $\dfrac{1}{3}+\dfrac{1}{6}+\dfrac{1}{9}$
    \item $\dfrac{5}{6}-\dfrac{1}{3}+\dfrac{1}{2}$
    \item $\dfrac{7}{8}-\dfrac{3}{4}+\dfrac{1}{2}$
    \item $\dfrac{2}{3}+\dfrac{1}{4}-\dfrac{5}{12}$
    \item $\dfrac{9}{10}-\left(\dfrac{1}{5}+\dfrac{1}{2}\right)$
  \end{subexercise}
\end{exercise}

\newpage

\chaptertitle{第五单元\quad 分数乘除}

\sectiontitle{方法回顾}
\knowbox{分数乘法：分子乘分子，分母乘分母，能约分先约分。\\
分数除法：除以一个数等于乘这个数的倒数。\\
带分数化成假分数后再运算。\\
乘法运算律对分数同样适用。}

\bigskip

\begin{exercise}
\item 计算：
  \begin{subexercise}
    \item $\dfrac{2}{3}\times\dfrac{3}{4}$
    \item $\dfrac{5}{6}\times\dfrac{2}{5}$
    \item $\dfrac{3}{8}\times\dfrac{4}{9}$
    \item $\dfrac{7}{12}\times\dfrac{6}{7}$
    \item $\dfrac{4}{15}\times\dfrac{5}{8}$
    \item $\dfrac{9}{20}\times\dfrac{10}{27}$
  \end{subexercise}

\item 计算：
  \begin{subexercise}
    \item $\dfrac{2}{3}\div\dfrac{4}{5}$
    \item $\dfrac{5}{6}\div\dfrac{5}{12}$
    \item $\dfrac{3}{8}\div\dfrac{9}{16}$
    \item $\dfrac{7}{10}\div\dfrac{14}{15}$
    \item $\dfrac{4}{9}\div\dfrac{8}{27}$
    \item $\dfrac{9}{14}\div\dfrac{3}{7}$
  \end{subexercise}

\item 计算：
  \begin{subexercise}
    \item $2\dfrac{1}{2}\times\dfrac{2}{5}$
    \item $3\dfrac{1}{3}\times\dfrac{3}{10}$
    \item $4\dfrac{2}{7}\div\dfrac{5}{7}$
    \item $2\dfrac{4}{5}\div\dfrac{7}{10}$
    \item $1\dfrac{3}{5}\times\dfrac{5}{8}$
    \item $3\dfrac{3}{4}\div\dfrac{5}{6}$
  \end{subexercise}

\item 简算：
  \begin{subexercise}
    \item $\dfrac{5}{7}\times\dfrac{3}{8}+\dfrac{5}{7}\times\dfrac{5}{8}$
    \item $\dfrac{4}{9}\times\dfrac{7}{11}+\dfrac{4}{9}\times\dfrac{4}{11}$
    \item $\dfrac{3}{8}\times\dfrac{5}{9}+\dfrac{3}{8}\times\dfrac{4}{9}$
    \item $\dfrac{8}{15}\times\dfrac{7}{13}+\dfrac{8}{15}\times\dfrac{6}{13}$
  \end{subexercise}

\item 计算：
  \begin{subexercise}
    \item $\left(\dfrac{1}{2}+\dfrac{1}{3}\right)\times6$
    \item $\left(\dfrac{2}{3}-\dfrac{1}{4}\right)\times12$
    \item $\left(\dfrac{3}{4}+\dfrac{5}{6}\right)\times24$
    \item $\left(\dfrac{7}{8}-\dfrac{5}{12}\right)\times48$
  \end{subexercise}

\item 计算：
  \begin{subexercise}
    \item $\dfrac{3}{5}\times\dfrac{5}{6}\times\dfrac{2}{3}$
    \item $\dfrac{4}{7}\times\dfrac{7}{8}\times\dfrac{2}{5}$
    \item $\dfrac{2}{3}\div\dfrac{4}{5}\times\dfrac{3}{5}$
    \item $\dfrac{5}{6}\times\dfrac{3}{10}\div\dfrac{1}{4}$
  \end{subexercise}
\end{exercise}

\newpage

\chaptertitle{第六单元\quad 分数小数四则混合}

\sectiontitle{方法回顾}
\knowbox{分数化小数：分子除以分母。\\
小数化分数：一位小数是十分之几，两位是百分之几\ldots\\
混合运算顺序：先乘除后加减，有括号先算括号。\\
能约分的先约分，能简算的要简算。}

\bigskip

\begin{exercise}
\item 将分数化成小数：
  \begin{subexercise}
    \item $\dfrac{3}{5}$
    \item $\dfrac{7}{8}$
    \item $\dfrac{2}{25}$
    \item $\dfrac{9}{20}$
    \item $\dfrac{7}{16}$
    \item $\dfrac{11}{40}$
  \end{subexercise}

\item 将小数化成分数并化简：
  \begin{subexercise}
    \item $0.6$
    \item $0.75$
    \item $0.125$
    \item $0.375$
    \item $0.65$
    \item $0.875$
  \end{subexercise}

\item 计算：
  \begin{subexercise}
    \item $\dfrac{1}{2}+0.3$
    \item $\dfrac{3}{4}-0.25$
    \item $0.6+\dfrac{2}{5}$
    \item $\dfrac{7}{8}-0.5$
    \item $1.2-\dfrac{3}{5}$
    \item $0.75+\dfrac{1}{3}$
  \end{subexercise}

\item 计算：
  \begin{subexercise}
    \item $\dfrac{2}{3}+0.5-0.25$
    \item $1.2-\dfrac{1}{4}+0.8$
    \item $2.5\times\dfrac{3}{5}$
    \item $\dfrac{4}{5}\times0.25$
    \item $3.6\div\dfrac{9}{10}$
    \item $\dfrac{5}{6}\div0.5$
  \end{subexercise}

\item 简算：
  \begin{subexercise}
    \item $0.25\times\dfrac{3}{7}\times4$
    \item $1.25\times\dfrac{5}{8}\times8$
    \item $4.8\times\dfrac{3}{8}+\dfrac{5}{8}\times4.8$
    \item $7.2\times\dfrac{3}{5}+2.8\times\dfrac{3}{5}$
  \end{subexercise}

\item 计算：
  \begin{subexercise}
    \item $\left(0.5+\dfrac{1}{3}\right)\times12$
    \item $\left(2.4-\dfrac{4}{5}\right)\div0.2$
    \item $3.2\times\left(\dfrac{3}{4}+0.5\right)$
    \item $\left(\dfrac{5}{6}-0.5\right)\div\dfrac{1}{3}$
  \end{subexercise}
\end{exercise}

\newpage

\chaptertitle{综合闯关一}

\begin{exercise}
\item $136+89+64+11$
\item $25\times17\times4\times3$
\item $199+299+399+99$
\item $125\times25\times8\times4$
\item $1000-345-255$
\item $36\times99+36$
\item $2.5\times3.6\times4$
\item $0.125\times88$
\item $4.8\div0.25+5.2\div0.25$
\item $\dfrac{2}{3}+\dfrac{1}{4}+\dfrac{1}{3}+\dfrac{3}{4}$
\item $\dfrac{5}{7}\times\dfrac{3}{8}+\dfrac{5}{7}\times\dfrac{5}{8}$
\item $\dfrac{3}{4}\times\dfrac{2}{5}\div\dfrac{3}{5}$
\item $\left(\dfrac{5}{6}-\dfrac{1}{3}\right)\times18$
\item $3.2\times\left(\dfrac{3}{8}+0.25\right)$
\item $\dfrac{5}{9}-\dfrac{2}{9}+\dfrac{4}{9}-\dfrac{1}{9}$
\end{exercise}

\newpage

\chaptertitle{综合闯关二}

\begin{exercise}
\item $9999+999+99+9$
\item $125\times72$
\item $25\times36$
\item $1999\times5$
\item $360\div45$
\item $450\div18$
\item $37\times101-37$
\item $4.5\times9.9+4.5\times0.1$
\item $12.5\times0.25\times32$
\item $\dfrac{1}{5}+\dfrac{1}{10}+\dfrac{1}{20}+\dfrac{1}{40}$
\item $\dfrac{7}{12}\times\dfrac{3}{4}+\dfrac{7}{12}\times\dfrac{1}{4}$
\item $\dfrac{8}{9}\times\dfrac{3}{4}\times\dfrac{3}{2}$
\item $1.2\div\dfrac{3}{5}\times0.5$
\item $\left(\dfrac{3}{4}-0.25\right)\div0.5$
\item $0.8\times\left(\dfrac{5}{8}+\dfrac{3}{8}\div0.5\right)$
\end{exercise}

\newpage

\chaptertitle{综合闯关三}

\begin{exercise}
\item $875+364+125+36$
\item $47\times63+37\times47$
\item $99\times88\div(9\times11)$
\item $5.6\times17.8-5.6\times7.8$
\item $0.25\times4.8\div0.2$
\item $36.5\times99+36.5$
\item $1.25\times88-1.25\times8$
\item $\dfrac{5}{6}+\dfrac{7}{12}-\dfrac{5}{6}+\dfrac{5}{12}$
\item $\left(\dfrac{5}{9}+\dfrac{3}{8}\right)\times72$
\item $\dfrac{4}{15}\times\dfrac{5}{8}+\dfrac{11}{15}\times\dfrac{5}{8}$
\item $3.5\times\dfrac{4}{7}\div0.5$
\item $\dfrac{2}{3}\times\left(0.75-\dfrac{1}{4}\right)$
\item $12.5\times0.25\times3.2$
\item $\dfrac{5}{8}\times0.4+\dfrac{3}{8}\div2.5$
\item $999\times999+1999$
\end{exercise}

\bigskip

\begin{center}
  \textbf{\large 全部答案见书末}
\end{center}

\newpage

\chaptertitle{答案}

\sectiontitle{第一单元}
\noindent
1.(a) 75 \quad (b) 157 \quad (c) 289 \quad (d) 647 \quad (e) 9999 \quad (f) 1999 \\
2.(a) 597 \quad (b) 1505 \quad (c) 2337 \quad (d) 11094 \quad (e) 1874 \quad (f) 5599 \\
3.(a) 250 \quad (b) 497 \quad (c) 1003 \quad (d) 1505 \\
4.(a) 536 \quad (b) 466 \quad (c) 693 \quad (d) 1000 \quad (e) 863 \quad (f) 1458 \\
5.(a) 726 \quad (b) 374 \quad (c) 532 \quad (d) 445 \quad (e) 1085 \quad (f) 3321 \\
6.(a) 556 \quad (b) 934 \quad (c) 745 \quad (d) 934 \quad (e) 832 \quad (f) 767 \\
7.(a) 144 \quad (b) 225 \quad (c) 211 \quad (d) 333 \quad (e) 500 \quad (f) 800

\sectiontitle{第二单元}
\noindent
1.(a) 1700 \quad (b) 23000 \quad (c) 3700 \quad (d) 1300 \quad (e) 7000 \quad (f) 900 \\
2.(a) 3564 \quad (b) 4848 \quad (c) 4975 \quad (d) 100125 \quad (e) 5508 \quad (f) 3626 \\
3.(a) 1100 \quad (b) 11000 \quad (c) 4410 \quad (d) 3136 \quad (e) 3600 \quad (f) 4700 \\
4.(a) 5600 \quad (b) 7200 \quad (c) 3800 \quad (d) 9500 \quad (e) 11000 \quad (f) 800 \\
5.(a) 16 \quad (b) 6.4 \quad (c) 24 \quad (d) 240 \quad (e) 4 \quad (f) 0.8 \\
6.(a) 10 \quad (b) 4 \quad (c) 8 \quad (d) 15 \quad (e) 15 \quad (f) 15 \\
7.(a) 1400 \quad (b) 800 \quad (c) 100000 \quad (d) 9900 \quad (e) 8787 \quad (f) 3400

\sectiontitle{第三单元}
\noindent
1.(a) 8.4 \quad (b) 3.3 \quad (c) 21.25 \quad (d) 15.76 \quad (e) 9.12 \quad (f) 3.42 \\
2.(a) 14 \quad (b) 3 \quad (c) 1.2 \quad (d) 1 \quad (e) 1.8 \quad (f) 0.3 \\
3.(a) 9 \quad (b) 30 \quad (c) 42 \quad (d) 2.7 \quad (e) 3 \quad (f) 16 \\
4.(a) 37 \quad (b) 7 \quad (c) 1.1 \quad (d) 1.1 \quad (e) 360 \quad (f) 480 \\
5.(a) 56 \quad (b) 45 \quad (c) 10 \quad (d) 101 \quad (e) 55.44 \quad (f) 34.34 \\
6.(a) 5.9 \quad (b) 16.4 \quad (c) 12 \quad (d) 6 \quad (e) 2.5 \quad (f) 10 \\
7.(a) 1.2 \quad (b) 0.5 \quad (c) 11.2 \quad (d) 6.5 \quad (e) 90 \quad (f) 2.7

\sectiontitle{第四单元}
\noindent
1.(a) $\frac35$ \quad (b) $\frac59$ \quad (c) 1 \quad (d) $\frac12$ \quad (e) $1\frac27$ \quad (f) $\frac13$ \\
2.(a) $\frac56$ \quad (b) $\frac{5}{12}$ \quad (c) $\frac{11}{12}$ \quad (d) $\frac{11}{24}$ \quad (e) $\frac{7}{10}$ \quad (f) $\frac14$ \\
3.(a) $\frac{29}{36}$ \quad (b) $\frac14$ \quad (c) $\frac{31}{72}$ \quad (d) $\frac{19}{90}$ \quad (e) $\frac{31}{42}$ \quad (f) $\frac{39}{140}$ \\
4.(a) $3\frac{7}{12}$ \quad (b) $2\frac{1}{15}$ \quad (c) $7\frac{13}{24}$ \quad (d) $2\frac16$ \quad (e) $9\frac{11}{14}$ \quad (f) $3\frac{1}{6}$ \\
5.(a) $1\frac25$ \quad (b) $2$ \quad (c) $1\frac{5}{11}$ \quad (d) $\frac{7}{10}$ \\
6.(a) $\frac78$ \quad (b) $\frac{11}{18}$ \quad (c) 1 \quad (d) $\frac58$ \quad (e) $\frac12$ \quad (f) 0

\sectiontitle{第五单元}
\noindent
1.(a) $\frac12$ \quad (b) $\frac13$ \quad (c) $\frac16$ \quad (d) $\frac12$ \quad (e) $\frac16$ \quad (f) $\frac{1}{6}$ \\
2.(a) $\frac56$ \quad (b) 2 \quad (c) $\frac23$ \quad (d) $\frac34$ \quad (e) $1\frac12$ \quad (f) $1\frac12$ \\
3.(a) 1 \quad (b) 1 \quad (c) 6 \quad (d) 4 \quad (e) 1 \quad (f) $4\frac12$ \\
4.(a) $\frac57$ \quad (b) $\frac49$ \quad (c) $\frac38$ \quad (d) $\frac{8}{15}$ \\
5.(a) 5 \quad (b) 5 \quad (c) 38 \quad (d) 22 \\
6.(a) $\frac13$ \quad (b) $\frac{4}{35}$ \quad (c) $\frac12$ \quad (d) $\frac14$

\sectiontitle{第六单元}
\noindent
1.(a) 0.6 \quad (b) 0.875 \quad (c) 0.08 \quad (d) 0.45 \quad (e) 0.4375 \quad (f) 0.275 \\
2.(a) $\frac35$ \quad (b) $\frac34$ \quad (c) $\frac18$ \quad (d) $\frac38$ \quad (e) $\frac{13}{20}$ \quad (f) $\frac78$ \\
3.(a) 0.8 \quad (b) 0.5 \quad (c) 1 \quad (d) 0.375 \quad (e) 0.6 \quad (f) $1\frac{1}{12}$ \\
4.(a) $0.91\frac23$ \quad (b) $1.75$ \quad (c) 1.5 \quad (d) 0.2 \quad (e) 4 \quad (f) $1\frac23$ \\
5.(a) $1\frac{2}{7}$ \quad (b) $7\frac{13}{16}$ \quad (c) 4.8 \quad (d) 6 \\
6.(a) 10 \quad (b) 8 \quad (c) 4 \quad (d) $1\frac{1}{3}$ \quad (省略单位)

\sectiontitle{综合闯关}
\noindent
闯关一：1. 300 \quad 2. 5100 \quad 3. 996 \quad 4. 100000 \quad 5. 400 \quad 6. 3600 \quad 7. 36 \quad 8. 11 \quad 9. 40 \quad 10. 2 \quad 11. $\frac57$ \quad 12. $\frac12$ \quad 13. 9 \quad 14. 2 \quad 15. $\frac29$ \\
闯关二：1. 11106 \quad 2. 9000 \quad 3. 900 \quad 4. 9995 \quad 5. 8 \quad 6. 25 \quad 7. 3700 \quad 8. 45 \quad 9. 100 \quad 10. $\frac{15}{40}$ \quad 11. $\frac{7}{12}$ \quad 12. 1 \quad 13. 1 \quad 14. 1 \quad 15. $1\frac14$ \\
闯关三：1. 1400 \quad 2. 4700 \quad 3. 88 \quad 4. 56 \quad 5. 6 \quad 6. 3650 \quad 7. 100 \quad 8. 1 \quad 9. 67 \quad 10. $\frac58$ \quad 11. 4 \quad 12. $\frac13$ \quad 13. 10 \quad 14. 0.4 \quad 15. 1000000

\end{document}
